\scanpage{17}
% archive foliation: 16
ces conditions, les entiers $\mu_i$ $(i\geq1)$,
$\nu_j$ $(j\geq0)$, $\beta_{i\ell}$ $(i\geq1,\ell\in\unsure{\mathfrak P})$
sont uniquement
déterminés par les conditions suivantes

a) Pour \inserted{tout} $i\geq1$, les $\beta_{i\ell}$ \inserted{$\neq0$}
sont les \inserted{multiplicités des} diviseurs
élémentaires de \deleted{\unsure{$H_i$}}
$\pi_{i-1}(L_*)$ $(=H_{i-1}(L^!_*))$

b) \deleted{\unsure{$k$}} $\nu_i=\operatorname{rg}\pi_i(L_*)$
$(=\operatorname{rg}H_i(L^!_*))$

c) On a \deleted{rang} $\mu_i=\operatorname{rg}B_{i-1}(L^!_*)$
\deleted{\unsure{$\beta_{i-1}$}} $-\beta_i$
\[
(\text{où }\beta_i=\sum_\ell\beta_{i,\ell})
\]

d) On a de plus
\[
\left\{\begin{aligned}
\operatorname{rg}Z_i\deleted{\unsure{\partial}}(L^!_*)
 &=\mu_{i+1}\deleted{\unsure{+\mu_i}}+\beta_{i+1}+\nu_i\\
\operatorname{rg}L^!_i
 &=\mu_{i+1}+\deleted{\unsure{\mu_i+}}\beta_{i+1}+\beta_i+\nu_i
\end{aligned}\right.
\]
Enfin $L_*$ est (str.) $\partial$-parfait ssi le nb de
facteurs non nuls qui interviennent est fini.

\uline{Dém} \deleted{\unsure{En fait}} \inserted{Par} la théorie de Dold-Puppe,
on réduit tout à une question de
complexes de chaînes. Comme
tt sous-module d'un $k$-module
projectif de t.f. sur $k$ est itou, on
conclut que tout tel complexe,
s'il est \deleted{\unsure{str. ps. coh.}} \inserted{str.} à comp. proj. de t.f.,
s'écrit comme un produit de
complexes (définis à isom \deleted{unique}
\deleted{\unsure{unique}} près, \deleted{\unsure{par le choix de}}
\inserted{--- mais l'isom}
\inserted{$\prod_i\simeq L^!_*$ n'est pas unique}) de la forme
\diagram{\begin{tikzcd}[column sep=large]
K_i:\quad -0\arrow[r] &
P_i\arrow[r,hook]\arrow[d,equal] &
Q_i\arrow[r]\arrow[d,equal] &0-\\
& B_{i-1}\arrow[r,phantom,"\subset"] & Z_{i-1} &
\end{tikzcd}}
\inserted{$P_i$: degré $i$\qquad $Q_i$: degré $i-1$}
